\renewcommand{\sectionauthor}{Prof. Gustavo Ale}

\begin{multicols*}{2}
        
    \section*{Geometria}
    A geometria é provavelmente a área mais antiga da matemática, precursora da própria álgebra e sendo um dos pilares 
    da matemática. A palavra geometria é resultado da combinação das palavras gregas \textit{geo} e 
    \textit{metron}/\textit{metri}, significando respectivamente Terra e medição, pois a área da geometria consiste da
    medição e entendimento das relações e propriedades  contidas nas fíguras geométricas: comprimento, distância, ângulo, 
    área, volume e perímetro, as fíguras geométricas por sua vez são parte integrante da natureza e da Terra (\textit{geo}).
    
    \subsection*{Definições gerais}
    Antes do estudo da geometria devemos ter ciência das definições de conceitos presentes nesse ramo da matemática, 
    dentro destes conceitos se enquadram as medidas, os ângulos, as formas geométricas e suas componentes.  
    
    \subsection*{Formas geométricas e suas componentes}
    A primeira componente das formas geométricas e a mais simples delas é o ponto, o mesmo deve ser imaginado como um 
    ponto infinitesimal, não possuindo comprimento, área ou perímetro. Quando um ponto é integrante de uma forma 
    geométrica o mais comum é chamarmos ele de \textit{vértice}

    \begin{figure}[H]
        \centering
        \includegraphics[width=0.7\columnwidth]{assets/image_placeholder.png}
        \caption{Vértice de um triângulo }
    \end{figure}

    A linha por sua vez é o elemento que liga dois pontos $P_1$ e $P_2$ em uma reta, assim como o ponto é idealmente 
    infinitesimal a linha idealmente não possui espessura. Quando uma linha é parte integrante de uma forma geometria 
    sua denominação comum é aresta.

    \begin{figure}[H]
        \centering
        \includegraphics[width=0.7\columnwidth]{assets/image_placeholder.png}
        \caption{Aresta de um quadrado}
    \end{figure}


    Quando se tem 3 ou mais pontos o resultado da conexão deles é tido como superfície, plano ou forma geométrica 
    dependendo da área da matemática e do contexto de estudo, porém vale ter em mente que também existem superfícies 
    não bidimensionais. No contexto da geometria a ser estudado nesse capítulo a conexão de 3 ou mais pontos 
    geram exclusivamente formas geométricas.

    \begin{figure}[H]
        \centering
        \includegraphics[width=0.7\columnwidth]{assets/image_placeholder.png}
        \caption{Um quadrilátero com as arestas realçadas em vermelho}
    \end{figure}




    % teste matemática em duas colunas
    \color{black}
    \begin{align}
        f(x) &= \frac{sen(x)}{cos(x)} \\[10pt]
        f(x) &= tan(x)
    \end{align}
    \color{mediumgray}

    %\pagebreak
    \newpage
    \section*{Ângulos}
    \lipsum[1]

    \subsection*{Simetrias de figuras planas}
    \lipsum[1-2]    
    
\end{multicols*}